% steps/step08.tex
\section[گام هشتم: ارزیابی مدل]{\faChartBar\ \ گام هشتم: ارزیابی مدل}
\begin{stepbox}

\field{چرا معیار صحت عمومی (Accuracy) فریبنده است؟}{
    در تصاویر CT Scan کبد، معمولاً تومورها بخش بسیار کوچکی (مثلاً کمتر از ۱٪) از کل حجم تصویر را اشغال می‌کنند. در چنین شرایطی با عدم تعادل شدید کلاس‌ها روبرو هستیم؛ به‌طوری که اگر مدل به‌طور ساده تمام پیکسل‌ها را پس‌زمینه پیش‌بینی کند، صحت عمومی سامانه ۹۹٪ ثبت خواهد شد، در حالی که مدل از لحاظ بالینی کاملاً بی‌ارزش است. به همین دلیل، ارزیابی مدل بر پایه معیارهای مبتنی بر هم‌پوشانی و سنجش فواصل مرزی انجام می‌گیرد.
}

\field{معیارهای ارزیابی مناسب چیست؟}{
    برای سنجش دقیق عملکرد مدل از شاخص‌های تخصصی زیر استفاده می‌شود:
    \begin{itemize}
        \item \textbf{ضریب دایس (Dice Score) و تقاطع روی اجتماع (IoU):} این دو معیار، میزان همپوشانی هندسی ماسک پیش‌بینی‌شده ($X$) و ماسک واقعی ($Y$) را ارزیابی می‌کنند. فرمول شاخص IoU به صورت زیر تعریف می‌شود:
        \[
            \mathrm{IoU} = \frac{|X \cap Y|}{|X \cup Y|}
        \]
        \item \textbf{فاصله هاسدورف (Hausdorff Distance - HD):} این معیار بیشینه فاصله میان مرزهای پیش‌بینی‌شده توسط مدل و مرزهای واقعی را بر حسب میلی‌متر محاسبه می‌کند. سنجش دقیق فاصله هاسدورف برای پزشکان و جراحان حیاتی است تا مطمئن شوند مدل مرز ضایعه را فراتر یا کمتر از حد واقعی تخمین نزده و در نتیجه در فرآیند درمان یا جراحی آسیبی به بافت‌های سالم مجاور وارد نمی‌شود.
    \end{itemize}
}

\field{تحلیل نتایج و شناسایی نقاط ضعف و قوت}{
    ارزیابی نهایی روی مجموعه آزمون (Test Set) که در فرآیند آموزش کاملاً دست‌نخورده باقی مانده است، انجام می‌گیرد. در این مرحله، عملکرد مدل در بخش‌بندی کبد سالم (کلاس آسان‌تر به دلیل ابعاد بزرگ‌تر و مرزهای واضح‌تر) و تومورها (کلاس دشوارتر به دلیل ابعاد متغیر و کنتراست پایین) به‌صورت مجزا تحلیل شده و چالش‌هایی نظیر خطای مدل در مواجهه با تومورهای بسیار کوچک یا موارد دارای آرتیفکت‌های شدید بررسی می‌گردند.
}

\end{stepbox}
